\chapter{Organização}

\section{Organograma}

O organograma foi desenhado a partir da página de liderança da AstraZeneca \cite{LEAD}. Mostra o Conselho de Administração na Tabela~\ref{tab:conselho} e o Chief Executive Officer (CEO), a Senior Executive Team (SET) e as unidades abaixo deles na Figura~\ref{fig:unidades}. A fonte de cada caixa é indicada na Tabela~\ref{tab:fontes-organograma}.

\begin{table}[htbp]\centering
\caption{Conselho de Administração: 14 administradores e cinco comissões (Auditoria, Remunerações, Nomeações e Governação, Ciência, Sustentabilidade) \cite[pp.~67--68]{AR}.}
\label{tab:conselho}
\begin{tabular}{@{}llc@{}}
\toprule
\textbf{Administrador} & \textbf{Cargo} & \textbf{No AR} \\
\midrule
Michel Demaré      & Presidente não executivo              & \textcolor{orgok}{sim} \\
Philip Broadley    & Administrador Independente Sénior     & \textcolor{orgok}{sim} \\
Euan Ashley        & Administrador não executivo           & \textcolor{orgok}{sim} \\
Birgit Conix       & Administradora não executiva          & \textcolor{orgok}{sim} \\
Karen Knudsen      & Administradora não executiva          & \textcolor{orgok}{sim} \\
Diana Layfield     & Administradora não executiva          & \textcolor{orgok}{sim} \\
Anna Manz          & Administradora não executiva          & \textcolor{orgok}{sim} \\
Sheri McCoy        & Administradora não executiva          & \textcolor{orgok}{sim} \\
Tony Mok           & Administrador não executivo           & \textcolor{orgok}{sim} \\
Marcus Wallenberg  & Administrador não executivo           & \textcolor{orgok}{sim} \\
Wan Ling Martello  & Apenas no organograma                 & \textcolor{orgdiff}{não} \\
\bottomrule
\end{tabular}
\end{table}

\begin{figure}[htbp]\centering
\resizebox{\linewidth}{!}{%
\begin{forest}
  forked edges,
  for tree={grow'=east, parent anchor=east, child anchor=west,
            edge={draw=azline, line width=0.6pt}, l sep=9pt, s sep=2.5pt},
  [{\nd{Pascal Soriot}{Presidente Executivo}}, ok
    [{\nd{Senior Executive Team}{10 membros}}, ok
      [{\nd{Financeira}{CFO: Aradhana Sarin}}, na]
      [{\nd{R\&D Therapeuticals}{agrupamento}}, grp
        [{\nd{Oncologia}{}}, grp
          [{\nd{D. Fredrickson}{EVP, Unidade de Negócio de Oncologia}}, dif]
          [{\nd{S. Galbraith}{EVP, I\&D de Oncologia}}, dif]
        ]
        [{\nd{BioPharma}{}}, grp
          [{\nd{S. Barr}{EVP, I\&D BioPharma}}, ok]
          [{\nd{R. Dobber}{EVP, BioPharma BU}}, ok]
        ]
        [{\nd{Doenças Raras}{}}, grp
          [{\nd{M. Dunoyer}{CEO Alexion \& Diretor de Estratégia}}, ok]
        ]
        [{\nd{Vacinas \& Terapias Imunitárias}{}}, dif
          [{\nd{R. Dobber}{também em BioPharma}}, ok]
        ]
      ]
      [{\nd{Comercial}{agrupamento}}, grp
        [{\nd{I. Reic}{EVP, Internacional}}, ok]
        [{\nd{Marketing Global}{}}, na]
        [{\nd{Reino Unido}{}}, na]
        [{\nd{Cone Sul}{}}, na]
      ]
      [{\nd{Operações}{agrupamento}}, grp
        [{\nd{P. Cheng}{EVP Operações Globais, IT \& CSO}}, ok]
        [{\nd{J. Pott}{CHRO, Compliance \& Diretor Jurídico}}, ok]
      ]
    ]
  ]
\end{forest}}

\medskip
{\sffamily\scriptsize
\textcolor{orgok}{\rule{7pt}{7pt}}~Corresponde ao AR \quad
\textcolor{orggrp}{\rule{7pt}{7pt}}~Agrupamento acrescentado para leitura \quad
\textcolor{orgdiff}{\rule{7pt}{7pt}}~Diverge do AR \quad
\textcolor{orgna}{\rule{7pt}{7pt}}~Não verificável no AR}

\caption{Do Chief Executive Officer até à Senior Executive Team e às unidades abaixo dela. Desenhado para este relatório a partir de \cite{LEAD}; a fonte de cada caixa é indicada na Tabela~\ref{tab:fontes-organograma}.}
\label{fig:unidades}
\end{figure}

\begin{xltabular}{\linewidth}{@{}>{\raggedright\arraybackslash}p{3.0cm}>{\raggedright\arraybackslash}p{2.9cm}Y@{}}
\caption{Fontes de cada caixa da Tabela~\ref{tab:conselho} e da Figura~\ref{fig:unidades}.}\label{tab:fontes-organograma}\\*
\toprule
\textbf{Caixa} & \textbf{Estado} & \textbf{Fonte e notas} \\*
\midrule
\endfirsthead
\toprule
\textbf{Caixa} & \textbf{Estado} & \textbf{Fonte e notas} \\*
\midrule
\endhead
\bottomrule
\endfoot
Conselho de Administração & Corresponde ao AR & O Conselho tem 14 administradores e cinco comissões (Auditoria, Remunerações, Nomeações e Governação, Ciência, Sustentabilidade). As matérias não reservadas ao Conselho são delegadas no CEO ou numa das comissões. \cite[pp.~67--68]{AR} \\
Michel Demaré & Corresponde ao AR & Presidente não executivo do Conselho. \cite[p.~68]{AR} \\
Philip Broadley & Corresponde ao AR & Administrador não executivo independente sénior. \cite[p.~68]{AR} \\
Administradores não executivos confirmados & Corresponde ao AR & Euan Ashley, Birgit Conix, Karen Knudsen, Diana Layfield, Anna Manz, Sheri McCoy, Tony Mok e Marcus Wallenberg constam do AR como administradores não executivos. \cite[pp.~68--69]{AR} \\
Administradores não executivos confirmados & Corresponde ao AR & Consta do AR como administrador não executivo. \cite[pp.~68--69]{AR} \\
Administradores não executivos confirmados & Corresponde ao AR & Consta do AR como administrador não executivo. \cite[pp.~68--69]{AR} \\
Administradores não executivos confirmados & Corresponde ao AR & Consta do AR como administrador não executivo. \cite[pp.~68--69]{AR} \\
Administradores não executivos confirmados & Corresponde ao AR & Consta do AR como administrador não executivo. \cite[pp.~68--69]{AR} \\
Administradores não executivos confirmados & Corresponde ao AR & Consta do AR como administrador não executivo. \cite[pp.~68--69]{AR} \\
Administradores não executivos confirmados & Corresponde ao AR & Consta do AR como administrador não executivo. \cite[pp.~68--69]{AR} \\
Marcus Wallenberg & Corresponde ao AR & Consta do AR como administrador não executivo. \cite[pp.~68--69]{AR} \\
Wan Ling Martello & Diverge do AR & Aparece neste organograma mas NÃO aparece na lista do Conselho do AR (verificado por pesquisa de texto). O AR lista ainda Rene Haas e Nazneen Rahman como administradores não executivos, que estão em falta no organograma. \cite[pp.~67--69]{AR} \\
Pascal Soriot, CEO & Corresponde ao AR & Administrador executivo e Chief Executive Officer; membro da SET. \cite[pp.~68, 70]{AR} \\
Senior Executive Team (SET) & Corresponde ao AR & A SET é o órgão através do qual o CEO exerce a autoridade delegada pelo Conselho. Tem 10 membros, cinco deles mulheres. \cite[pp.~39, 70]{AR} \\
Caixa ``Financeira'' & Não verificável no AR & Caixa sem titular neste organograma. A CFO, Aradhana Sarin, é administradora executiva e membro da SET \cite[pp.~68, 70]{AR}, e no organograma está ligada ao Conselho e à SET. \\
``R\&D Therapeuticals'' (caixa de agrupamento) & Agrupamento acrescentado para leitura & Caixa de agrupamento acrescentada para facilitar a leitura. O AR descreve três organizações de I\&D por área terapêutica (Oncologia, BioPharmaceuticals, Rare Disease) \cite[p.~26]{AR} mas não apresenta esta caixa como uma unidade. \\
Oncologia & Agrupamento acrescentado para leitura & O AR lista a Oncologia como uma das três organizações de I\&D \cite[p.~26]{AR}. No organograma, os dois responsáveis de Oncologia ficam sob esta caixa. \\
David Fredrickson & Diverge do AR & Título no AR: EVP, Oncology Haematology Business Unit. O organograma omite ``Haematology''. \cite[p.~70]{AR} \\
Susan Galbraith & Diverge do AR & Título no AR: EVP, Oncology Haematology R\&D. O organograma omite ``Haematology''. \cite[p.~70]{AR} \\
BioPharma & Agrupamento acrescentado para leitura & A BioPharmaceuticals é uma das três organizações de I\&D. \cite[p.~26]{AR} \\
Sharon Barr & Corresponde ao AR & EVP, BioPharmaceuticals R\&D. \cite[p.~70]{AR} \\
Ruud Dobber & Corresponde ao AR & EVP, BioPharmaceuticals Business Unit, responsável por CVRM, R\&I e Vacinas \& Terapias Imunitárias. Isto explica por que razão o organograma o liga às duas caixas. Uma só pessoa a cobrir duas áreas não é reporte dual. \cite[p.~70]{AR} \\
Doenças Raras & Agrupamento acrescentado para leitura & A Rare Disease é uma das três organizações de I\&D. \cite[p.~26]{AR} \\
Marc Dunoyer & Corresponde ao AR & CEO da Alexion e Diretor de Estratégia da AstraZeneca, membro da SET. Dois cargos oficiais ocupados por uma só pessoa. \cite[p.~70]{AR} \\
Vacinas \& Terapias Imunitárias & Diverge do AR & No AR, isto faz parte da BioPharmaceuticals (área de Ruud Dobber) e não é uma organização de I\&D separada. \cite[pp.~26, 70]{AR} \\
Ruud Dobber (segunda ligação) & Corresponde ao AR & A mesma pessoa que aparece em BioPharma. \cite[p.~70]{AR} \\
``Comercial'' (caixa de agrupamento) & Agrupamento acrescentado para leitura & Caixa de agrupamento acrescentada para facilitar a leitura. O AR descreve três regiões comerciais: EUA, Europa-Canadá e Internacional; o Japão reporta em separado. Só a Internacional aparece no organograma. \cite[p.~32]{AR} \\
Iskra Reic & Corresponde ao AR & EVP, Internacional: China, mercados asiáticos e da Eurásia, Médio Oriente \& África, América Latina, Austrália e Nova Zelândia. \cite[p.~70]{AR} \\
Marketing Global & Não verificável no AR & Não foi encontrado no AR e é mantido apenas como agrupamento do organograma. \\
Reino Unido & Não verificável no AR & Não é apresentado como unidade no AR. \\
Cone Sul & Não verificável no AR & Não aparece no AR. \\
Operações (caixa de agrupamento) & Agrupamento acrescentado para leitura & Caixa de agrupamento acrescentada para facilitar a leitura. No AR, Pam Cheng e Jeff Pott são membros distintos da SET. \cite[p.~70]{AR} \\
Pam Cheng & Corresponde ao AR & EVP, Operações Globais, IT e Chief Sustainability Officer. \cite[p.~70]{AR} \\
Jeff Pott & Corresponde ao AR & Chief Human Resources Officer, Chief Compliance Officer e Diretor Jurídico (o AR lista também a propriedade intelectual). \cite[p.~70]{AR} \\
\end{xltabular}

O Relatório e Contas confirma a cadeia Conselho $\rightarrow$ CEO $\rightarrow$ SET, os dez membros da SET e os respetivos títulos \cite[p.~70]{AR}, e a existência de três organizações de I\&D por área terapêutica e de três regiões comerciais \cite[pp.~26, 32]{AR}. As caixas ``R\&D Therapeuticals'', ``Comercial'', ``Operações'' e ``Financeira'' não são unidades do Grupo: agrupam o organograma para que possa ser lido de uma só vez. A análise que segue assenta nas unidades que o Relatório e Contas documenta.

\section{Análise da estrutura}

Os critérios abaixo vêm do material da unidade curricular sobre a função organização \cite{ORG}. Cada afirmação limita-se ao que o AR, o organograma ou o material da unidade curricular sustentam, e cada análise termina com aquilo que não se pode concluir a partir das fontes.

\subsection{Tipo de departamentalização}

Departamentalização significa agrupar atividades em unidades lógicas para executar a estratégia \cite[p.~14]{ORG}. O material da unidade curricular distingue o agrupamento funcional \cite[p.~15]{ORG}, o agrupamento divisional por produto, área geográfica, cliente ou projeto \cite[pp.~16--19]{ORG}, o agrupamento múltiplo (misto), quando as grandes organizações não podem apoiar-se num único critério \cite[p.~20]{ORG}, e as estruturas matriciais, que combinam dois critérios e têm autoridade dual, com os trabalhadores a reportar a dois gestores \cite[pp.~21--22]{ORG}.

A AstraZeneca usa vários critérios ao mesmo tempo:

\begin{itemize}
\item \textbf{Por área terapêutica (produto):} três organizações de I\&D --- Oncologia, BioPharmaceuticals e Rare Disease \cite[p.~26]{AR} ---, com um EVP distinto para a I\&D e para a unidade de negócio na Oncologia e na BioPharmaceuticals \cite[p.~70]{AR}.
\item \textbf{Por área geográfica:} três regiões comerciais (EUA, Europa-Canadá, Internacional), com o Japão a reportar em separado \cite[p.~32]{AR}.
\item \textbf{Por função:} as Operações (fabrico e abastecimento, IT), a área financeira e os recursos humanos, compliance e jurídico são liderados por membros distintos da SET \cite[p.~70]{AR}.
\end{itemize}

\begin{analysis}
A estrutura descreve-se melhor como \textbf{departamentalização múltipla (mista)}, combinando os critérios funcional, por área terapêutica e geográfico \cite[p.~20]{ORG}.
\end{analysis}

\begin{sourceslimit}
Não se pode chamar-lhe matricial. Uma matriz exige autoridade dual \cite[pp.~21--22]{ORG}, e nem o organograma nem o AR documentam alguém a reportar a dois gestores. A única ligação dupla do organograma --- Ruud Dobber em BioPharma e em Vacinas \& Terapias Imunitárias --- é uma só pessoa a cobrir duas áreas \cite[p.~70]{AR}. O AR diz também que o Grupo tem um único segmento reportável e que a descoberta, o desenvolvimento, o fabrico e a comercialização são geridos globalmente numa base altamente integrada \cite[p.~147]{AR}.
\end{sourceslimit}

\subsection{Capacidade de adaptação à mudança}

O material da unidade curricular liga a função organização à gestão da mudança \cite[pp.~4--5]{ORG}, diz que as estruturas com poucos níveis e com delegação respondem mais depressa do que as que têm muitos níveis e centralização \cite[p.~7]{ORG}, e que as estruturas funcionais são rígidas e difíceis de reestruturar \cite[p.~16]{ORG}.

O AR documenta uma mudança deliberada e prolongada no tempo:

\begin{itemize}
\item Depois de adquirir a Alexion em 2021, o Grupo lançou o Post Alexion Acquisition Group Review, um programa global de reestruturação para integrar sistemas, estrutura e processos. Custou \$237m em 2025 e \$3,414m em termos acumulados \cite[p.~142]{AR}, com um custo total previsto de \$4.0bn \cite[p.~59]{AR}.
\item ``Enable an agile organisation'' é uma das áreas de foco estratégico, ao lado de ``Transform R\&D ways of working'' \cite[p.~11]{AR}.
\item Em dezembro de 2025, o Grupo criou uma unidade dedicada de Enterprise AI para acelerar a Ambition 2030 \cite[p.~36]{AR}.
\item Marc Dunoyer lidera a Alexion dentro da SET \cite[p.~70]{AR}, e Sharon Barr e a CFO, Aradhana Sarin, trabalharam anteriormente na Alexion \cite[pp.~68, 70]{AR}.
\end{itemize}

\begin{analysis}
O Grupo mostrou que consegue absorver uma aquisição de grande dimensão e reorganizar-se de propósito, e aponta a agilidade como objetivo. Isto indica uma \textbf{capacidade de adaptação razoável}.
\end{analysis}

\begin{sourceslimit}
O AR não dá tempos de decisão, pelo que não é possível avaliar a \textit{rapidez} da resposta. Os percursos de carreira de executivos individuais (por exemplo, Iskra Reic) e uma ``reorganização de 2024'' não constam do AR, pelo que não são usados.
\end{sourceslimit}

\subsection{Organização formal e informal}

A organização formal é concebida pela gestão, tem regras e procedimentos e pode ser desenhada num organograma. A organização informal surge da interação social, não está nem pode estar no organograma, e pode acelerar a comunicação ou resistir à mudança \cite[pp.~12--13]{ORG}.

\textbf{Formal:} o Conselho reserva para si certas matérias e delega as restantes no CEO ou nas suas cinco comissões \cite[p.~67]{AR}; a SET é o órgão através do qual o CEO exerce a autoridade delegada \cite[p.~70]{AR}; existe um Código de Ética e um programa global de compliance \cite[p.~34]{AR}; e subcomissões transversais e comités de produto alocam recursos \cite[p.~147]{AR}. O facto de Marc Dunoyer ocupar dois cargos --- CEO da Alexion e Diretor de Estratégia --- é também \textit{formal}, porque ambos são posições oficiais \cite[p.~70]{AR}.

\textbf{Ligação documentada mais próxima da organização informal:} os Employee Resource Groups, que são grupos voluntários, conduzidos pelos trabalhadores e abertos a todos os trabalhadores \cite[p.~39]{AR}, e o inquérito Pulse aos trabalhadores \cite[p.~39]{AR}.

\begin{sourceslimit}
Nenhum documento descreve a organização informal real (redes, comunidades de prática, líderes informais). Por definição, o organograma não a mostra \cite[p.~12]{ORG}. O relatório explica, por isso, apenas o conceito e os canais organizados que existem.
\end{sourceslimit}

\subsection{Flexibilidade ou rigidez}

Segundo o material da unidade curricular, as grandes organizações tendem a ser mais especializadas, mais centralizadas e mais formalizadas \cite[p.~11]{ORG}, e a rigidez cresce com a formalização \cite[p.~9]{ORG} e com a centralização \cite[p.~7]{ORG}.

\begin{itemize}
\item Todas as decisões operacionais significativas são tomadas ao nível da SET, e os recursos são alocados a todo o Grupo por comités de produto \cite[p.~147]{AR}. Isto indica \textbf{centralização} no topo.
\item A governação está documentada em detalhe: matérias reservadas, delegações, comissões e respetivos papéis \cite[pp.~67, 71]{AR}.
\item A presença comercial está organizada através de empresas de marketing locais detidas na totalidade na maioria dos mercados \cite[p.~32]{AR}, e o Grupo declara um foco numa organização ágil \cite[p.~11]{AR}.
\end{itemize}

\begin{analysis}
A evidência aponta para uma estrutura \textbf{formalizada e centralizada na SET}, com especialização por área terapêutica e por região, e com a agilidade como objetivo declarado. Tende para a rigidez na tomada de decisão e para a flexibilidade apenas na intenção declarada.
\end{analysis}

\begin{sourceslimit}
O AR não descreve autonomia para as unidades abaixo da SET, nem autonomia regional de preços. A ideia de que a Alexion funciona como uma ``empresa autónoma'' não é sustentada: o AR diz que o Grupo tem um único segmento reportável e que as áreas funcionais não são geridas em separado \cite[p.~147]{AR}.
\end{sourceslimit}

\subsection{Número de níveis hierárquicos}

O material da unidade curricular relaciona os níveis com a amplitude de controlo: quanto maior a amplitude, menos níveis e mais achatada a estrutura; amplitudes menores dão mais níveis e uma estrutura alta \cite[pp.~8--9]{ORG}. Nomeia também três níveis de gestão: estratégico, tático e operacional \cite[p.~7]{ORG}.

\begin{table}[htbp]\centering
\caption{Níveis contados na Figura~\ref{fig:unidades}, incluindo o Conselho como nível~1.}
\label{tab:niveis}
\begin{tabularx}{\linewidth}{@{}llY@{}}
\toprule
\textbf{Ramo} & \textbf{Níveis} & \textbf{Cadeia} \\
\midrule
Operações & 5 & Conselho $\rightarrow$ CEO $\rightarrow$ SET $\rightarrow$ Operações $\rightarrow$ executivo (4 abaixo do Conselho) \\
\addlinespace
I\&D e Comercial & 6 & Um nível a mais, por causa das caixas de agrupamento R\&D Therapeuticals e Comercial (5 abaixo do Conselho) \\
\bottomrule
\end{tabularx}
\end{table}
\begin{sourceslimit}
A contagem refere-se ao \textbf{organograma analisado}, que cobre apenas o topo da empresa. O AR não divulga os níveis abaixo da SET nem as amplitudes de controlo, pelo que não se pode concluir se a estrutura no seu todo é alta ou achatada. Outras contagens de níveis que aparecem no ficheiro de pesquisa (6 a 9 níveis) não têm fonte no AR e não são usadas.
\end{sourceslimit}

\subsection{Posições de linha e de staff}

Um gestor de \textbf{linha} tem poder de decisão e autoridade sobre os trabalhadores que lhe reportam, dentro da cadeia de comando. Uma unidade ou pessoa de \textbf{staff} apoia, aconselha ou assessora os gestores de linha. No exemplo do material da unidade curricular, os diretores de marketing, financeiro e de produção são posições de linha, enquanto o secretariado da direção e a assessoria jurídica são staff \cite[p.~8]{ORG}.

\begin{xltabular}{\linewidth}{@{}YYY@{}}
\caption{Posições de linha e de staff na Senior Executive Team. Classificação: análise dos autores com base em \cite[p.~8]{ORG}.}\label{tab:linha-staff}\\*
\toprule
\textbf{Posição} & \textbf{Responsabilidade segundo o AR} & \textbf{Classificação} \\*
\midrule
\endfirsthead
\toprule
\textbf{Posição} & \textbf{Responsabilidade segundo o AR} & \textbf{Classificação} \\*
\midrule
\endhead
\bottomrule
\endfoot
EVP da I\&D e das unidades de negócio: Fredrickson, Galbraith, Barr, Dobber & Conduzem a I\&D e o desempenho comercial das suas áreas terapêuticas \cite[p.~70]{AR} & \textbf{Linha} \\
Iskra Reic, EVP Internacional & Estratégia e crescimento na região Internacional \cite[p.~70]{AR} & \textbf{Linha} \\
Pam Cheng, EVP Operações Globais, IT e CSO & Operações, IT e sustentabilidade \cite[p.~70]{AR} & \textbf{Mista} --- o fabrico é linha; o IT e a sustentabilidade apoiam outras áreas \\
Jeff Pott, CHRO, Compliance e Diretor Jurídico & Recursos humanos, compliance, jurídico e propriedade intelectual \cite[p.~70]{AR} & \textbf{Staff} --- a assessoria jurídica é staff no exemplo da unidade curricular \\
Aradhana Sarin (CFO); Marc Dunoyer (Alexion, Diretor de Estratégia) & Área financeira; Alexion e estratégia \cite[p.~70]{AR} & \textbf{Não classificada} --- o AR não diz quem lhes reporta \\
\end{xltabular}

\section{Propostas de alteração à estrutura}

Cada proposta abaixo assenta num facto do AR e nos compromissos descritos no material da unidade curricular. São propostas, não alterações anunciadas.

\subsection{Proposta 1 · Coordenação mais forte entre a I\&D e a unidade de negócio em cada área terapêutica}

\textbf{Base.} Na Oncologia e na BioPharmaceuticals, a I\&D e a unidade de negócio são lideradas por membros diferentes da SET \cite[p.~70]{AR}. O AR afirma que já existem subcomissões transversais e comités de produto \cite[p.~147]{AR}.

\textbf{Proposta.} Tornar a coordenação ao nível da área terapêutica permanente e visível, com um fórum conjunto por área onde a I\&D e a unidade de negócio partilham prioridades e recursos.

\textbf{Vantagem.} Decisões mais rápidas e menos lacunas entre a investigação e os planos comerciais. \textbf{Risco.} Mais reuniões e decisões mais lentas se o fórum não tiver autoridade clara. O material da unidade curricular avisa que, com unidades baseadas no produto, os interesses de cada unidade podem passar à frente dos da organização e as unidades podem competir em vez de cooperar \cite[p.~17]{ORG}; um fórum conjunto serve para contrariar isso.

\subsection{Proposta 2 · Mais autonomia local no acesso e nos preços}

\textbf{Base.} A pressão sobre os preços por parte dos pagadores é uma das principais preocupações comerciais em várias regiões \cite[p.~32]{AR}. Hoje, o Grupo aplica uma política global de preços por escalões, com quatro escalões baseados no rendimento nacional bruto \cite[p.~41]{AR}, e as decisões operacionais são tomadas ao nível da SET \cite[p.~147]{AR}.

\textbf{Proposta.} Dar às equipas comerciais regionais mais margem para adaptar o acesso e os preços, dentro de limites definidos centralmente pela SET.

\textbf{Vantagem.} A delegação aproxima as decisões do mercado e acelera a resposta \cite[p.~7]{ORG}, e, quando as unidades estão muito distantes entre si, é mais adequado atribuir a gestão a um gestor localizado na região \cite[p.~17]{ORG}. \textbf{Risco.} Menos consistência entre países; os limites têm de ser claros. O AR não reporta quanta autonomia as regiões têm hoje, pelo que esta proposta necessita dessa informação para ser testada.
